\section{Validation}
\label{sec:validation}

We validate \policyengine{} UK against three independent sources, each
targeting a different level of aggregation: population-level aggregates
against an independent microsimulation model and official administrative/
forecast figures (Section~\ref{sec:validation-ukmod}); individual-level tax
liabilities against \hmrc{} administrative microdata
(Section~\ref{sec:validation-spi}); and a further individual-level check on
a specific, harder-to-model program (Section~\ref{sec:validation-other}).
Every figure quoted in this section is drawn from a committed,
version-controlled source in \texttt{policyengine-uk} or \texttt{populace};
none is recomputed for this paper. Where a number needs to be refreshed
against the model's current state before submission, it is marked
\todo{example}.

\subsection{Aggregate validation: \ukmod{} and administrative outturns}
\label{sec:validation-ukmod}

\texttt{docs/book/validation/validation.ipynb} in \texttt{policyengine-uk}
runs a three-way comparison -- \policyengine{}, \ukmod{}, and an official
target -- for major tax and benefit programs across 2023--2025: income tax,
National Insurance (employee, self-employed, and employer contributions
combined), Universal Credit, Child Benefit, Tax Credits, Housing Benefit,
and Pension Credit. \ukmod{} figures are transcribed from the ``2020-26
UKMOD country report'' (cited in-notebook; \todo{resolve to a full, citable
reference -- ISER/EUROMOD country report series, exact edition and year}).
The ``Official'' series is \policyengine{}'s own calibration target,
sourced from \obr{}/\hmrc{}/\dwp{} budgetary-impact parameters (e.g.
\texttt{programs.income\_tax.budgetary\_impact.by\_country.UNITED\_KINGDOM}
for income tax; the equivalent \texttt{total\_NI}, \texttt{universal\_credit},
\texttt{child\_benefit}, \texttt{tax\_credits}, \texttt{housing\_benefit}, and
\texttt{pension\_credit} parameters for the other programs), i.e. it
represents the administrative/forecast aggregate the model is calibrated
to, not an independent third-party figure -- this distinction should be
stated plainly in the drafted prose so the comparison is not misread as
\policyengine{} vs.\ two independent externals.

This notebook has a maintained history: it was most recently updated for a
new \ukmod{} country report in policyengine-uk\#724 (``Update validation
page for new UKMOD country report''), and earlier issues
(policyengine-uk\#300, ``List differences against UKMOD in documentation'';
\#212, ``Validate employer NI reform against UKMOD report''; \#545,
``Generate a UKMOD input dataset'') document the comparison's development
history. \todo{Re-run the notebook against the current model version before
the paper drafts numeric claims -- the notebook's own import of
\texttt{policyengine\_uk.data.datasets.frs.calibration.loss} no longer
resolves against the current \texttt{policyengine-uk} tree (calibration
code has moved into \texttt{populace.calibrate}), so it must be updated to
run before its output can be cited with specific figures in prose. This
outline cites the notebook's existence, method, and program coverage, which
are independently verifiable by reading its source; it does not assert
specific match percentages or error magnitudes that would require rerunning
it.}

Public framing constraint (per issue scope): this comparison is descriptive
and validation-comparative only. It reports where \policyengine{}, \ukmod{},
and the official target agree or disagree on point estimates, with no
characterization of which model is ``more correct'' beyond distance from
the cited official figure, and no commentary on \ukmod{} or \euromod{}
beyond their own published numbers.

\subsection{Individual-level validation: the Survey of Personal Incomes}
\label{sec:validation-spi}

\texttt{docs/book/validation/spi-validation.ipynb} validates
\policyengine{} UK's calculated individual income tax liability against
\hmrc{}'s \spi{} 2020/21 microdata, using \spi{} input variables as direct
inputs to household-level \policyengine{} simulations (excluding the
\spi{}'s 1,770 ``composite records,'' about 0.2\% of the file, which
represent population aggregates rather than individuals and cannot be
compared this way). Matches are logarithmically banded by absolute error,
with values within \pounds{}10 treated as matches (attributed to rounding).
The notebook reports:

\begin{table}[H]
    \centering
    \caption{PolicyEngine UK vs.\ SPI 2020/21: income tax match rate
    (source: \texttt{docs/book/validation/spi-validation.ipynb})}
    \label{tab:spi-match}
    \tablefont
    \begin{tabular}{lc}
        \toprule
        Total error & Share of records \\
        \midrule
        $\leq$\pounds{}10 & 92.7\% \\
        $<$\pounds{}100 & 93.5\% \\
        $<$\pounds{}1{,}000 & 98.3\% \\
        $\geq$\pounds{}1{,}000 & 1.7\% \\
        \bottomrule
    \end{tabular}
    \tablenote{Cumulative bands as reported in the notebook (i.e. rows are
    not mutually exclusive: ``$<$\pounds{}100'' includes the
    ``$\leq$\pounds{}10'' records). \todo{Confirm cumulative vs.\ exclusive
    banding before reproducing this table in prose -- the source markdown
    table's rows sum to more than 100\% read literally, consistent with a
    cumulative reading, but this should be verified against the underlying
    code rather than assumed.}}
\end{table}

The notebook attributes the residual 6.5\% (records with error $\geq
\pounds{}100$) to four identified discrepancy clusters, each with an
estimated record share: Marriage Allowance transfer bookkeeping (up to
4.8\% of records), non-resident Personal Allowance treatment (0.3\%), an
``Other Investment Income'' savings-income tax discrepancy (0.2\%), and
Gift Aid interaction with the Personal Allowance taper (0.1\%). This
validation is referenced in the notebook itself as building on an earlier
public writeup, \policyengine{}'s ``Autumn 2023 model calibration''
(\url{https://policyengine.org/uk/research/uk-calibration-2023-q4},
confirmed live). \todo{This validation uses \spi{} 2020/21 -- the most
recent vintage available at time of writing per the notebook's own citation
to the \spi{} statistics release. Confirm whether a newer \spi{} vintage has
since been released and, if so, whether to note that this validation
predates it or to rerun against the newer vintage before submission.}

\subsection{Other individual-level checks: student loan repayments}
\label{sec:validation-other}

\texttt{docs/book/validation/student-loan-repayments.ipynb} compares
\frs{}-\emph{reported} student loan repayments against \policyengine{}'s
\emph{modelled} repayments (calculated from income and loan-plan type under
the Plan 1/2/4/5 and Postgraduate repayment thresholds). Unlike the \spi{}
comparison, this is a within-\frs{} check (reported vs.\ modelled on the
same records) rather than a cross-dataset one. The notebook finds aggregate
totals ``reasonably aligned'' but individual-level correlation
``lower than for income tax calculations,'' attributed to five factors:
within-year employment/income variation not captured by an annual model,
multiple concurrent loan plans, current-student repayment patterns, and
loan-plan misclassification in the imputed microdata. \todo{Extract and cite
the specific correlation coefficient, aggregate ratio, and match-rate
figures the notebook computes (it defines \texttt{correlation},
\texttt{total\_reported}, \texttt{total\_modelled}, and \texttt{match\_rate}
variables) -- rerun the notebook to get current values rather than
transcribing stale ones, since (like Section~\ref{sec:validation-ukmod}) its
numeric outputs are not cached in this outline.}

This is presented as a validation of a harder, more idiosyncratic program
rather than a headline result: the paper should be explicit that this
program's individual-level fit is weaker than income tax's, consistent with
the model's own documentation of that finding, rather than presenting only
favorable comparisons.

\subsection{Validation not yet assembled}

\todo{The issue scope also names ``HMRC/DWP ready reckoners'' as a target
validation source, distinct from the budgetary-impact calibration
parameters already covered in Section~\ref{sec:validation-ukmod}. A ready
reckoner (e.g. HMRC's ``Direct effects of illustrative tax changes'' or
DWP's benefit expenditure tables) gives the marginal fiscal effect of a
specific parameter change, which is a distinct and currently-unassembled
validation target: comparing \policyengine{}'s reform-costing output against
a published ready-reckoner line item for the same parameter change. No such
comparison was found in the policyengine-uk or populace repositories during
this survey (2026-07-04); this is a genuine evidence gap, not an
oversight in this outline, and should be either assembled before v1
circulates or explicitly scoped out with a stated reason.}

\todo{HBAI-level distributional validation (poverty rates, income deciles)
against DWP's own published HBAI statistics is a natural complement to the
program-level aggregate validation in Section~\ref{sec:validation-ukmod} --
the model's HBAI income-concept mapping (\texttt{docs/book/validation/hbai.md})
establishes that the necessary variables exist, but no notebook comparing
\policyengine{}-calculated HBAI poverty/income statistics against DWP's
published HBAI statistics was found. State whether this is in scope for v1
or deferred.}
